\SourceLeaf{356}%
gering, dass man es bei den meisten Messungen kaum für die geringe Mühe wert halten wird, sie zu berücksichtigen.

In der Abhandlung wird die Theorie aller dieser und weiterer damit zusammenhängender Rechnungen vollständig entwickelt und an einer durchgehenden Musterrechnung erläutert. Eine Hilfstafel ist beigefügt, die zwar zunächst für die Zone bestimmt ist, in der das hannoversche Dreiecksnetz liegt, sich aber ohne Weiteres auch für Vermessungen verwenden lässt, die weit über diese Zone hinausreichen. Sie erstreckt sich nämlich über eine Zone von zwölf Breitengraden, in deren Mitte der gewählte Normalparallelkreis auf \(52^0\,40'\) Breite liegt. Die Tafel ist so genau berechnet, dass sie selbst ausreicht, wenn ein Dreiecksnetz mit zehnstelligen Logarithmen berechnet werden soll. Ihre Genauigkeit ist also weit höher als diejenige, die man in den meisten Fällen beibehalten wird. Die geringe Platzersparnis, die das Weglassen einiger Dezimalstellen gebracht hätte, schien jedoch zu unbedeutend, um beim Abdruck etwas davon wegzulassen.

Merklich und unmerklich sind bei Rechnungen relative Begriffe. Es lohnt sich daher, sie anhand von ein paar Beispielen aus der Abhandlung auf ein bestimmtes Maß zurückzuführen.

Im hannoverschen Dreiecksnetz ist das größte Dreieck, das zugleich am weitesten vom Normalparallelkreis entfernt liegt, das zwischen den Punkten Brocken, Hohenhagen und Inselsberg gebildete. In ihm kommen daher auch die größten Werte der Richtungsreduktionen vor, und zwar an der Seite Hohenhagen-Inselsberg, wo die Reduktion des Azimuts am ersten Endpunkt \(-0''00332\) und am anderen \(+0''00428\) beträgt. Im ganzen Netz gibt es nur noch zwei weitere Dreiecksseiten, bei denen die Reduktionen \(0''001\) überschreiten. Bei allen übrigen bleiben sie unter diesem Wert.

Das größte Hauptdreieck der trigonometrischen Vermessungen der Schweiz ist das zwischen den Punkten Chasseral, Suchet und Berra gelegene. Es berührt gerade die Südgrenze, bis zu der sich die Hilfstafel erstreckt, sodass sich die Richtungsreduktionen noch mit ihrer Hilfe berechnen lassen. Die größte Reduktion betrifft das Azimut von Chasseral in Suchet und beträgt \(+0''06221\).

Daraus ist ersichtlich, dass in der ganzen Zone, in der das hannoversche Dreiecksnetz liegt, die Reduktion ganz entfallen kann, wenn die Rechnung auf
